\chapter{Experimental Setup and Results}
\label{chap:experiments}

\section{Experimental Setup}

Experiments are conducted on the Yamanishi08 dataset under the warm-start 1:1
and 1:10 settings described in Chapter~\ref{chap:data}. DistMult and CompGCN
are evaluated as alternative knowledge graph encoders within the same KGE-NFM
framework. The downstream feature construction, NFM architecture, data splits,
and evaluation procedure are kept consistent to support a controlled encoder
comparison.

\begin{table}[H]
\centering
\caption{Experimental configuration}
\label{tab:experiment_setup}
\begin{tabularx}{\textwidth}{lX}
\toprule
\textbf{Item} & \textbf{Setting} \\
\midrule
Dataset & Yamanishi08 \\
Evaluation settings & \texttt{warm\_start\_1\_1} and \texttt{warm\_start\_1\_10} \\
Number of supplied splits & 10 per setting \\
KGE encoders & DistMult and CompGCN \\
Embedding dimension & 400 \\
KGE training epochs & 50 \\
CompGCN layers & 1 \\
CompGCN composition & Element-wise multiplication \\
CompGCN dropout & 0.0 \\
Inverse triples & Enabled for CompGCN \\
NFM training epochs & Up to 2,000 with early stopping \\
Random seed & 0 \\
Device & Automatically selected CPU or CUDA device \\
\bottomrule
\end{tabularx}
\end{table}

With one CompGCN layer, the entity and relation updates in
Equations~\ref{eq:compgcn_node_update} and
\ref{eq:compgcn_relation_update} are applied once before embedding extraction.
The multiplicative composition
$\phi(\mathbf{h}_{u},\mathbf{h}_{r})=
\mathbf{h}_{u}\odot\mathbf{h}_{r}$ makes each neighbor message
relation-dependent, while the learned original, inverse, and self-loop
transformations control its direction. DistMult does not perform this
neighborhood update, so the controlled comparison isolates the effect of
relational graph convolution.

\section{Evaluation Metrics}

\subsection{ROC-AUC}

ROC-AUC measures how well the model ranks positive interactions above negative
or unknown interactions across all classification thresholds.

\subsection{PR-AUC}

PR-AUC summarizes the precision-recall curve. This metric is especially useful
for DTI prediction because positive interactions are usually sparse.

\section{Baselines}

% TODO: Add the actual baselines evaluated in this project.

\begin{table}[H]
\centering
\caption{Baseline models}
\label{tab:baseline_models}
\begin{tabularx}{\textwidth}{lX}
\toprule
\textbf{Model} & \textbf{Description} \\
\midrule
Descriptor-only classifier & Uses drug and protein descriptors without KGE features \\
Standalone KGE score & Uses only KGE interaction scores for prediction \\
KGE + Random Forest & Uses KGE features with a random forest classifier \\
KGE + XGBoost & Uses KGE features with an XGBoost classifier \\
KGE-NFM & Framework combining KG embeddings, descriptors, and NFM \\
\bottomrule
\end{tabularx}
\end{table}

\section{Results}

This section reports the mean and standard deviation across the ten supplied
warm-start splits. ROC-AUC measures overall ranking quality, whereas PR-AUC is
emphasized when interpreting the imbalanced 1:10 setting. All placeholders in
this section should be replaced using the final experiment outputs.

\subsection{Overall Performance Comparison}

Table~\ref{tab:main_results} presents the principal comparison among the
descriptor-only model and the two KGE-NFM encoder configurations. Reporting
both evaluation settings in one table makes it possible to compare model
quality and sensitivity to class imbalance.

\begin{table}[H]
\centering
\caption{Mean DTI prediction performance across the supplied splits}
\label{tab:main_results}
\small
\begin{tabular}{lcccc}
\toprule
& \multicolumn{2}{c}{\textbf{Warm-start 1:1}} &
\multicolumn{2}{c}{\textbf{Warm-start 1:10}} \\
\cmidrule(lr){2-3}\cmidrule(lr){4-5}
\textbf{Model} & \textbf{ROC-AUC} & \textbf{PR-AUC} &
\textbf{ROC-AUC} & \textbf{PR-AUC} \\
\midrule
Descriptor-only NFM & $-- \pm --$ & $-- \pm --$ & $-- \pm --$ & $-- \pm --$ \\
DistMult KGE-NFM & $-- \pm --$ & $-- \pm --$ & $-- \pm --$ & $-- \pm --$ \\
CompGCN KGE-NFM & $-- \pm --$ & $-- \pm --$ & $-- \pm --$ & $-- \pm --$ \\
\bottomrule
\end{tabular}
\end{table}

% Replace the framed box with:
% \includegraphics[width=0.9\textwidth]{images/results/main_performance_comparison.png}
\begin{figure}[H]
    \centering
    \fbox{\parbox[c][4.0cm][c]{0.86\textwidth}{\centering
    Result figure placeholder\\[4pt]
    Grouped ROC-AUC and PR-AUC comparison for both warm-start settings}}
    \caption{Performance comparison of the evaluated models under warm-start
    1:1 and 1:10.}
    \label{fig:main_performance_comparison}
\end{figure}

The final discussion should identify the best-performing model in each setting
and quantify its improvement over the descriptor-only baseline. Improvements
should be reported as absolute metric differences rather than only as relative
percentages.

\subsection{Results under Warm-Start 1:1}

The 1:1 setting contains approximately equal numbers of positive and sampled
negative pairs. This experiment measures model discrimination under balanced
class exposure. Table~\ref{tab:warm_start_1_1_results} should contain the
per-model mean and standard deviation, together with the best and worst split
scores to show result stability.

\begin{table}[H]
\centering
\caption{Detailed results under the warm-start 1:1 setting}
\label{tab:warm_start_1_1_results}
\small
\begin{tabular}{lrrrr}
\toprule
\textbf{Model} & \textbf{ROC-AUC} & \textbf{PR-AUC} &
\textbf{Best PR-AUC} & \textbf{Worst PR-AUC} \\
\midrule
Descriptor-only NFM & $-- \pm --$ & $-- \pm --$ & -- & -- \\
DistMult KGE-NFM & $-- \pm --$ & $-- \pm --$ & -- & -- \\
CompGCN KGE-NFM & $-- \pm --$ & $-- \pm --$ & -- & -- \\
\bottomrule
\end{tabular}
\end{table}

% TODO: State whether graph representations improve over descriptor-only NFM
% and whether the improvement is consistent across the ten supplied splits.

\subsection{Results under Warm-Start 1:10}

The 1:10 setting is more demanding because each positive pair is evaluated
against approximately ten sampled negatives. PR-AUC is the primary metric in
this setting because it directly reflects the model's ability to preserve
precision when positive interactions are rare.

\begin{table}[H]
\centering
\caption{Detailed results under the warm-start 1:10 setting}
\label{tab:warm_start_1_10_results}
\small
\begin{tabular}{lrrrr}
\toprule
\textbf{Model} & \textbf{ROC-AUC} & \textbf{PR-AUC} &
\textbf{Best PR-AUC} & \textbf{Worst PR-AUC} \\
\midrule
Descriptor-only NFM & $-- \pm --$ & $-- \pm --$ & -- & -- \\
DistMult KGE-NFM & $-- \pm --$ & $-- \pm --$ & -- & -- \\
CompGCN KGE-NFM & $-- \pm --$ & $-- \pm --$ & -- & -- \\
\bottomrule
\end{tabular}
\end{table}

% Replace the framed box with:
% \includegraphics[width=0.82\textwidth]{images/results/pr_curves_1_10.png}
\begin{figure}[H]
    \centering
    \fbox{\parbox[c][4.0cm][c]{0.78\textwidth}{\centering
    Result figure placeholder\\[4pt]
    Mean precision--recall curves for warm-start 1:10}}
    \caption{Mean precision--recall curves under the warm-start 1:10 setting.}
    \label{fig:pr_curves_1_10}
\end{figure}

% TODO: Compare each model with the random PR-AUC reference determined by the
% positive prevalence, and explain any reduction relative to the 1:1 setting.

\subsection{DistMult and CompGCN Encoder Comparison}

To isolate the effect of the knowledge graph encoder, DistMult and CompGCN are
compared using the same data splits, descriptor features, NFM architecture, and
training procedure. For split $i$, the paired improvement is defined as

\begin{equation}
    \Delta_i = M_{\text{CompGCN},i} - M_{\text{DistMult},i},
    \label{eq:paired_encoder_improvement}
\end{equation}

where $M$ denotes ROC-AUC or PR-AUC. A positive value indicates that CompGCN
outperforms DistMult on the same split.

\begin{table}[H]
\centering
\caption{Paired CompGCN improvement over DistMult}
\label{tab:encoder_paired_comparison}
\small
\begin{tabular}{lrrr}
\toprule
\textbf{Setting and metric} & \textbf{Mean $\Delta$} &
\textbf{Improved splits} & \textbf{$p$-value} \\
\midrule
Warm-start 1:1 ROC-AUC & -- & -- / 10 & -- \\
Warm-start 1:1 PR-AUC & -- & -- / 10 & -- \\
Warm-start 1:10 ROC-AUC & -- & -- / 10 & -- \\
Warm-start 1:10 PR-AUC & -- & -- / 10 & -- \\
\bottomrule
\end{tabular}
\end{table}

% Replace the framed box with:
% \includegraphics[width=0.85\textwidth]{images/results/encoder_paired_differences.png}
\begin{figure}[H]
    \centering
    \fbox{\parbox[c][4.0cm][c]{0.80\textwidth}{\centering
    Result figure placeholder\\[4pt]
    Per-split CompGCN minus DistMult metric differences}}
    \caption{Paired performance differences between CompGCN and DistMult.}
    \label{fig:encoder_paired_differences}
\end{figure}

The final text should state whether CompGCN improves performance consistently
or only on selected splits. If a paired significance test is used, its name and
interpretation should be reported without treating a non-significant result as
evidence that the models are equivalent.

\subsection{Effect of Class Imbalance}

This analysis compares each model's performance between the 1:1 and 1:10
settings. Since the positive interactions are unchanged, the difference mainly
reflects the larger number of sampled negative candidates. The discussion
should emphasize changes in PR-AUC and false-positive control; ROC-AUC alone may
hide degradation under increased imbalance.

\begin{table}[H]
\centering
\caption{Performance change from warm-start 1:1 to 1:10}
\label{tab:imbalance_effect}
\small
\begin{tabular}{lrr}
\toprule
\textbf{Model} & \textbf{$\Delta$ ROC-AUC} & \textbf{$\Delta$ PR-AUC} \\
\midrule
Descriptor-only NFM & -- & -- \\
DistMult KGE-NFM & -- & -- \\
CompGCN KGE-NFM & -- & -- \\
\bottomrule
\end{tabular}
\end{table}

% TODO: Identify which model is least sensitive to the change in class ratio
% and relate the observation to the scarcity analysis in Chapter 3.

\section{Ablation Study}

An ablation study helps measure the contribution of each feature source or
model component.

\begin{table}[H]
\centering
\caption{Ablation study template}
\label{tab:ablation_study}
\begin{tabular}{lcc}
\toprule
\textbf{Configuration} & \textbf{ROC-AUC} & \textbf{PR-AUC} \\
\midrule
Descriptors only & -- & -- \\
KGE embeddings only & -- & -- \\
Descriptors + KGE embeddings & -- & -- \\
Full KGE-NFM pipeline & -- & -- \\
\bottomrule
\end{tabular}
\end{table}

\section{Discussion}

% TODO: Discuss the final results. Suggested points:
% - Does KGE improve over descriptor-only baselines?
% - Does NFM improve over classical classifiers?
% - Are ROC-AUC and PR-AUC consistent?
% - Which split is hardest and why?

\section{Error Analysis}

% TODO: Add analysis of false positives and false negatives. Consider whether
% errors are concentrated around rare drugs, rare proteins, missing KG entities,
% or cold-start cases.

\section{Limitations}

The current project has several expected limitations:

\begin{itemize}
    \item negative labels may include unknown positive interactions;
    \item the current evaluation is limited to warm-start settings and does not
    measure generalization to entirely unseen drugs or targets;
    \item hyperparameter search may be limited by available computation;
    \item model predictions require biological validation before practical use.
\end{itemize}
